% Options for packages loaded elsewhere
% Options for packages loaded elsewhere
\PassOptionsToPackage{unicode}{hyperref}
\PassOptionsToPackage{hyphens}{url}
%
\documentclass[
  10pt,
  ignorenonframetext,
  aspectratio=169,
  brazilian,
]{beamer}
\newif\ifbibliography
\usepackage{pgfpages}
\setbeamertemplate{caption}[numbered]
\setbeamertemplate{caption label separator}{: }
\setbeamercolor{caption name}{fg=normal text.fg}
\beamertemplatenavigationsymbolsempty
% remove section numbering
\setbeamertemplate{part page}{
  \centering
  \begin{beamercolorbox}[sep=16pt,center]{part title}
    \usebeamerfont{part title}\insertpart\par
  \end{beamercolorbox}
}
\setbeamertemplate{section page}{
  \centering
  \begin{beamercolorbox}[sep=12pt,center]{section title}
    \usebeamerfont{section title}\insertsection\par
  \end{beamercolorbox}
}
\setbeamertemplate{subsection page}{
  \centering
  \begin{beamercolorbox}[sep=8pt,center]{subsection title}
    \usebeamerfont{subsection title}\insertsubsection\par
  \end{beamercolorbox}
}
% Prevent slide breaks in the middle of a paragraph
\widowpenalties 1 10000
\raggedbottom
\AtBeginPart{
  \frame{\partpage}
}
\AtBeginSection{
  \ifbibliography
  \else
    \frame{\sectionpage}
  \fi
}
\AtBeginSubsection{
  \frame{\subsectionpage}
}
\usepackage{iftex}
\ifPDFTeX
  \usepackage[T1]{fontenc}
  \usepackage[utf8]{inputenc}
  \usepackage{textcomp} % provide euro and other symbols
\else % if luatex or xetex
  \usepackage{unicode-math} % this also loads fontspec
  \defaultfontfeatures{Scale=MatchLowercase}
  \defaultfontfeatures[\rmfamily]{Ligatures=TeX,Scale=1}
\fi
\usepackage{lmodern}

\ifPDFTeX\else
  % xetex/luatex font selection
\fi
% Use upquote if available, for straight quotes in verbatim environments
\IfFileExists{upquote.sty}{\usepackage{upquote}}{}
\IfFileExists{microtype.sty}{% use microtype if available
  \usepackage[]{microtype}
  \UseMicrotypeSet[protrusion]{basicmath} % disable protrusion for tt fonts
}{}
\makeatletter
\@ifundefined{KOMAClassName}{% if non-KOMA class
  \IfFileExists{parskip.sty}{%
    \usepackage{parskip}
  }{% else
    \setlength{\parindent}{0pt}
    \setlength{\parskip}{6pt plus 2pt minus 1pt}}
}{% if KOMA class
  \KOMAoptions{parskip=half}}
\makeatother


\usepackage{longtable,booktabs,array}
\newcounter{none} % for unnumbered tables
\usepackage{calc} % for calculating minipage widths
\usepackage{caption}
% Make caption package work with longtable
\makeatletter
\def\fnum@table{\tablename~\thetable}
\makeatother
\usepackage{graphicx}
\makeatletter
\newsavebox\pandoc@box
\newcommand*\pandocbounded[1]{% scales image to fit in text height/width
  \sbox\pandoc@box{#1}%
  \Gscale@div\@tempa{\textheight}{\dimexpr\ht\pandoc@box+\dp\pandoc@box\relax}%
  \Gscale@div\@tempb{\linewidth}{\wd\pandoc@box}%
  \ifdim\@tempb\p@<\@tempa\p@\let\@tempa\@tempb\fi% select the smaller of both
  \ifdim\@tempa\p@<\p@\scalebox{\@tempa}{\usebox\pandoc@box}%
  \else\usebox{\pandoc@box}%
  \fi%
}
% Set default figure placement to htbp
\def\fps@figure{htbp}
\makeatother



\ifLuaTeX
\usepackage[bidi=basic,shorthands=off]{babel}
\else
\usepackage[bidi=default,shorthands=off]{babel}
\fi
\ifLuaTeX
  \usepackage{selnolig} % disable illegal ligatures
\fi


\setlength{\emergencystretch}{3em} % prevent overfull lines

\providecommand{\tightlist}{%
  \setlength{\itemsep}{0pt}\setlength{\parskip}{0pt}}



 


\makeatletter
\@ifpackageloaded{caption}{}{\usepackage{caption}}
\AtBeginDocument{%
\ifdefined\contentsname
  \renewcommand*\contentsname{Índice}
\else
  \newcommand\contentsname{Índice}
\fi
\ifdefined\listfigurename
  \renewcommand*\listfigurename{Lista de Figuras}
\else
  \newcommand\listfigurename{Lista de Figuras}
\fi
\ifdefined\listtablename
  \renewcommand*\listtablename{Lista de Tabelas}
\else
  \newcommand\listtablename{Lista de Tabelas}
\fi
\ifdefined\figurename
  \renewcommand*\figurename{Figura}
\else
  \newcommand\figurename{Figura}
\fi
\ifdefined\tablename
  \renewcommand*\tablename{Tabela}
\else
  \newcommand\tablename{Tabela}
\fi
}
\@ifpackageloaded{float}{}{\usepackage{float}}
\floatstyle{ruled}
\@ifundefined{c@chapter}{\newfloat{codelisting}{h}{lop}}{\newfloat{codelisting}{h}{lop}[chapter]}
\floatname{codelisting}{Listagem}
\newcommand*\listoflistings{\listof{codelisting}{Lista de Listagens}}
\makeatother
\makeatletter
\makeatother
\makeatletter
\@ifpackageloaded{caption}{}{\usepackage{caption}}
\@ifpackageloaded{subcaption}{}{\usepackage{subcaption}}
\makeatother

\usepackage{bookmark}
\IfFileExists{xurl.sty}{\usepackage{xurl}}{} % add URL line breaks if available
\urlstyle{same}
% fallback for those not using the hyperref driver hyperxmp:
\makeatletter
\@ifundefined{xmpquote}{\newcommand{\xmpquote}[1]{#1}}{}
\makeatother
\hypersetup{
  pdftitle={Parte I : Aprendizado Supervisionado},
  pdflang={pt-BR},
  hidelinks,
  pdfcreator={LaTeX via pandoc}}


\title{\texorpdfstring{Parte I : Aprendizado
Supervisionado}{Parte I : Aprendizado Supervisionado}}
\subtitle{\texorpdfstring{{Aprendizado em Alta
Dimensão}}{Aprendizado em Alta Dimensão}}
\author{}
\date{}

\begin{document}
\frame{\titlepage}


\begin{frame}{Onde estamos?}
\protect\phantomsection\label{onde-estamos}
Estudamos diferentes formas de estimar

\[
r(\mathbf x)=\mathbb E(Y\mid\mathbf X=\mathbf x)
\]

\begin{itemize}
\tightlist
\item
  KNN e regressão local utilizam \textbf{vizinhanças}.
\item
  Splines utilizam \textbf{funções de base}.
\item
  Árvores e ensembles utilizam \textbf{partições}.
\item
  Redes neurais aprendem \textbf{composições de transformações}.
\end{itemize}

\textbf{Como esses métodos se comportam quando o número de variáveis
aumenta?}

\note{Abertura e motivação: 10 minutos, incluindo os quatro slides
seguintes. Conectar a aula à diversidade de estimadores já vista na
disciplina.}
\end{frame}

\begin{frame}{Dois estudos com \(n=200\) e \(p=1.000\)}
\protect\phantomsection\label{dois-estudos-com-n200-e-p1.000}
\begin{columns}[T]
\begin{column}{0.5\linewidth}
\begin{block}{Estudo A: sensores}
\protect\phantomsection\label{estudo-a-sensores}
1.000 medições refletem principalmente \textbf{três condições de
funcionamento}.

Muitas variáveis compartilham informação.

\textbf{Hipótese: redundância.}
\end{block}
\end{column}

\begin{column}{0.5\linewidth}
\begin{block}{Estudo B: características moleculares}
\protect\phantomsection\label{estudo-b-caracteruxedsticas-moleculares}
A resposta depende de \textbf{cinco entre 1.000 variáveis}.

A maioria das variáveis é desnecessária para a predição, dadas as
relevantes.

\textbf{Hipótese: esparsidade.}
\end{block}
\end{column}
\end{columns}

\textbf{Os estudos têm a mesma dificuldade apenas porque possuem os
mesmos \(n\) e \(p\)?}
\end{frame}

\begin{frame}{Dimensões diferentes}
\protect\phantomsection\label{dimensuxf5es-diferentes}
{\def\LTcaptype{none} % do not increment counter
\begin{longtable}[]{@{}ll@{}}
\toprule\noalign{}
Quantidade & Significado \\
\midrule\noalign{}
\endhead
\(n\) & Número de observações \\
\(p\) & Número de variáveis observadas \\
\(d\) & Dimensão de uma representação dos dados \\
\(s\) & Número de variáveis ativas na função de regressão \\
\(m\) & Número de parâmetros do modelo \\
\bottomrule\noalign{}
\end{longtable}
}

Podemos ter \(d\ll p\), \(s\ll p\) e, em uma RNA, \(m\gg p\).

\textbf{Essas quantidades descrevem aspectos diferentes do problema.}

\note{A dimensão intrínseca da distribuição de X e a dimensão suficiente
para prever Y podem diferir. Nesta aula, d denota a dimensão da
representação indicada em cada exemplo.}
\end{frame}

\section{Dimensionalidade e tamanho
amostral}\label{dimensionalidade-e-tamanho-amostral}

\begin{frame}{Uma região pequena em muitas dimensões}
\protect\phantomsection\label{uma-regiuxe3o-pequena-em-muitas-dimensuxf5es}
Considere entradas uniformes em \([0,1]^p\).

Uma região com comprimento \(h\) em cada coordenada ocupa a fração

\[
\boxed{V=h^p}
\]

do volume total. Para \(h=0{,}2\):

\[
p=1:\ 0{,}2,\qquad
p=2:\ 0{,}04,\qquad
p=5:\ 0{,}00032.
\]

\textbf{A mesma resolução por coordenada cobre uma fração cada vez menor
dos dados.}

\note{Bloco sobre dimensionalidade: cerca de 20 minutos. O hipercubo
unitário tem volume total 1 em qualquer dimensão. O que decresce é a
fração ocupada por uma região com lado fixo h menor que 1. Não dizer que
o volume do hipercubo unitário cresce exponencialmente.}
\end{frame}

\begin{frame}{Quantas observações caem nessa região?}
\protect\phantomsection\label{quantas-observauxe7uxf5es-caem-nessa-regiuxe3o}
Se \(N_R\) conta as observações na região:

\[
N_R\sim\operatorname{Binomial}(n,h^p),
\qquad \mathbb E(N_R)=nh^p.
\]

Para obter \textbf{\(k\) observações em média}:

\[
\boxed{n=\frac{k}{h^p}.}
\]

Com \(k=20\) e \(h=0{,}2\), temos \(n=20\times5^p\).

Isso ilustra cobertura do espaço sob uniformidade, sem garantir uma
contagem mínima.
\end{frame}

\begin{frame}{A amostra cresce rapidamente}
\protect\phantomsection\label{a-amostra-cresce-rapidamente}
{\def\LTcaptype{none} % do not increment counter
\begin{longtable}[]{@{}rrr@{}}
\toprule\noalign{}
Dimensão \(p\) & Fração \(0{,}2^p\) & \(n\) para 20 observações em
média \\
\midrule\noalign{}
\endhead
1 & \(0{,}2\) & 100 \\
2 & \(0{,}04\) & 500 \\
3 & \(0{,}008\) & 2.500 \\
5 & \(0{,}00032\) & 62.500 \\
10 & \(0{,}0000001024\) & 195.312.500 \\
\bottomrule\noalign{}
\end{longtable}
}

\textbf{Esta conta não constitui uma recomendação amostral para qualquer
modelo.}
\end{frame}

\begin{frame}{Dimensão e cobertura do espaço}
\protect\phantomsection\label{dimensuxe3o-e-cobertura-do-espauxe7o}
\pandocbounded{\includegraphics[keepaspectratio]{aula_outros_aspectos_files/figure-beamer/unnamed-chunk-2-1.png}}

Manter o mesmo tamanho de vizinhança em cada coordenada exige
crescimento exponencial de \(n\).

\note{Cada unidade no eixo vertical representa multiplicar n por 10.
Gráfico calculado da expressão geométrica, sem ajuste ou simulação.}
\end{frame}

\begin{frame}{O que acontece com o KNN?}
\protect\phantomsection\label{o-que-acontece-com-o-knn}
\textbf{Situação:} queremos estimar o preço de uma casa usando imóveis
semelhantes. Com duas características, é fácil encontrar vizinhos
próximos; com dezenas, até o imóvel ``mais próximo'' pode ser muito
diferente.

Para encontrar \(k\) vizinhos com poucos dados:

\begin{itemize}
\tightlist
\item
  Precisamos ampliar a região de busca.
\item
  Os vizinhos podem diferir em várias coordenadas.
\item
  Variáveis irrelevantes podem dominar a distância.
\end{itemize}

\[
d^2(\mathbf x,\mathbf x')=
\underbrace{\sum_{j\in S}(x_j-x'_j)^2}_{\text{variáveis relevantes}}
+
\underbrace{\sum_{j\notin S}(x_j-x'_j)^2}_{\text{demais variáveis}}.
\]

Padronizar evita dominância de escala, mas mantém as variáveis
irrelevantes na distância.
\end{frame}

\begin{frame}{Regressão local e splines multivariadas}
\protect\phantomsection\label{regressuxe3o-local-e-splines-multivariadas}
\textbf{Situação:} em regressão linear local, ajustamos uma reta usando
observações próximas do ponto de interesse. Em muitas dimensões,
precisamos escolher entre uma vizinhança quase vazia e uma vizinhança
tão ampla que deixa de ser local.

\textbf{Regressão local:} vizinhanças maiores podem aumentar o viés,
enquanto regiões menores contêm poucos dados.

\textbf{Bases tensoriais:} com \(K\) funções por variável, uma
construção completa usa

\[
\boxed{K^p\text{ funções de base}.}
\]

Com \(K=5\) e \(p=10\): \(5^{10}=9.765.625\) funções.

\textbf{Restrições como aditividade reduzem essa complexidade.}
\end{frame}

\begin{frame}{Uma taxa que resume a dificuldade}
\protect\phantomsection\label{uma-taxa-que-resume-a-dificuldade}
Imagine tentar prever o consumo de energia usando temperatura, umidade e
outras condições ambientais. Se a relação for desconhecida e bastante
flexível, precisamos observar combinações dessas variáveis. A taxa
abaixo descreve como essa dificuldade cresce com o número de entradas.

Para uma classe de funções com suavidade \(\alpha\), sob condições
usuais:

\[
\text{erro de estimação integrado}\asymp
n^{-\frac{2\alpha}{2\alpha+p}}.
\]

Com \(\alpha=2\):

{\def\LTcaptype{none} % do not increment counter
\begin{longtable}[]{@{}rr@{}}
\toprule\noalign{}
Dimensão & Taxa \\
\midrule\noalign{}
\endhead
\(p=1\) & \(n^{-4/5}\) \\
\(p=4\) & \(n^{-1/2}\) \\
\(p=20\) & \(n^{-1/6}\) \\
\bottomrule\noalign{}
\end{longtable}
}

O crescimento de \(p\) reduz a velocidade de aprendizagem.

\note{Taxa minimax típica para erro quadrático integrado sobre classes
suaves, com hipóteses sobre suporte, desenho e ruído. Refere-se ao
excesso de erro em relação à função verdadeira, não inclui o ruído
irredutível de uma nova resposta. Fonte: Wasserman, notas Nonparametric
Regression, referência ao final.}
\end{frame}

\section{Redundância}\label{redunduxe2ncia}

\begin{frame}{Muitas medições, poucos fatores}
\protect\phantomsection\label{muitas-mediuxe7uxf5es-poucos-fatores}
\textbf{Situação:} uma indústria registra temperatura em 100 sensores.
Eles podem medir, com ruídos diferentes, uma mesma condição térmica. Há
100 colunas, mas talvez poucas fontes independentes de informação.

Suponha que as variáveis obedeçam a

\[
X_j=a_jZ+\eta_j,\qquad j=1,\ldots,p.
\]

\begin{itemize}
\tightlist
\item
  \(Z\): fator compartilhado.
\item
  \(a_j\): intensidade e direção da associação com o fator.
\item
  \(\eta_j\): componente específico da medição.
\end{itemize}

\textbf{Muitas variáveis podem medir aspectos de um mesmo fenômeno.}

\note{Bloco sobre redundância: cerca de 15 minutos. Exemplo: dezenas de
sensores medindo um estado térmico comum. Com ruído em todas as
coordenadas, a distribuição pode ter suporte de dimensão p, ainda que a
estrutura de sinal seja aproximadamente de baixa dimensão.}
\end{frame}

\begin{frame}{Redundância em um exemplo simulado}
\protect\phantomsection\label{redunduxe2ncia-em-um-exemplo-simulado}
\pandocbounded{\includegraphics[keepaspectratio]{aula_outros_aspectos_files/figure-beamer/unnamed-chunk-3-1.png}}

As duas coordenadas contêm informação fortemente compartilhada.
\end{frame}

\begin{frame}{Representação de baixa dimensão}
\protect\phantomsection\label{representauxe7uxe3o-de-baixa-dimensuxe3o}
Uma generalização é

\[
\mathbf X=A\mathbf Z+\boldsymbol\eta,
\qquad \mathbf Z\in\mathbb R^d,\quad d\ll p.
\]

Se a resposta depende principalmente desses fatores,

\[
Y=g(\mathbf Z)+\varepsilon,
\]

uma representação de dimensão \(d\) pode simplificar a predição.

\textbf{Precisamos preservar o sinal relacionado à resposta.}
\end{frame}

\begin{frame}{Variância das entradas e sinal preditivo}
\protect\phantomsection\label{variuxe2ncia-das-entradas-e-sinal-preditivo}
Considere variáveis independentes:

\[
X_1\sim N(0,100),\qquad X_2\sim N(0,1),
\]

\[
Y=5X_2+\varepsilon.
\]

\textbf{Sem padronização}, a primeira componente populacional acompanha
\(X_1\).

Entretanto, a resposta depende de \(X_2\).

\textbf{Explicar a variação de \(\mathbf X\) não garante preservar a
informação sobre \(Y\).}

\note{Padronizar altera este exemplo, mas a ACP continua sem utilizar Y.
Evitar concluir que toda componente de alta variância é irrelevante.
Fonte conceitual: notas de redução dimensional STATS 202, Stanford.}
\end{frame}

\begin{frame}{Métodos para explorar redundância}
\protect\phantomsection\label{muxe9todos-para-explorar-redunduxe2ncia}
{\def\LTcaptype{none} % do not increment counter
\begin{longtable}[]{@{}
  >{\raggedright\arraybackslash}p{(\linewidth - 4\tabcolsep) * \real{0.3333}}
  >{\raggedright\arraybackslash}p{(\linewidth - 4\tabcolsep) * \real{0.3333}}
  >{\raggedright\arraybackslash}p{(\linewidth - 4\tabcolsep) * \real{0.3333}}@{}}
\toprule\noalign{}
\begin{minipage}[b]{\linewidth}\raggedright
Método
\end{minipage} & \begin{minipage}[b]{\linewidth}\raggedright
Estratégia
\end{minipage} & \begin{minipage}[b]{\linewidth}\raggedright
Cuidado
\end{minipage} \\
\midrule\noalign{}
\endhead
Ridge & Contrai coeficientes e estabiliza o ajuste & Mantém todas as
variáveis \\
Modelo fatorial & Representa fatores compartilhados & Depende da
adequação dos fatores \\
\bottomrule\noalign{}
\end{longtable}
}

Com redundância, selecionar uma única variável de cada grupo pode perder
informação útil das demais medições.
\end{frame}

\section{Esparsidade}\label{esparsidade}

\begin{frame}{Poucas variáveis determinam a resposta}
\protect\phantomsection\label{poucas-variuxe1veis-determinam-a-resposta}
Considere \(p=1.000\), mas

\[
Y=3X_2-2X_{17}+1{,}5X_{204}+\varepsilon.
\]

A função de regressão utiliza apenas \(s=3\) variáveis.

\[
\boxed{s\ll p.}
\]

\textbf{O desafio inclui descobrir quais variáveis são relevantes.}

\note{Bloco sobre esparsidade: cerca de 15 minutos. Aqui relevância
significa participação na função condicional adotada no modelo. Uma
variável fora do suporte ainda pode ter associação marginal com Y por
correlação.}
\end{frame}

\begin{frame}{Esparsidade dos coeficientes}
\protect\phantomsection\label{esparsidade-dos-coeficientes}
\textbf{Situação:} um laboratório mede mil características, mas a
resposta depende de apenas cinco. A dificuldade não é apenas estimar os
cinco efeitos: é descobrir quais são eles.

\pandocbounded{\includegraphics[keepaspectratio]{aula_outros_aspectos_files/figure-beamer/unnamed-chunk-4-1.png}}

Ilustração com 50 variáveis: apenas três efeitos ativos.
\end{frame}

\begin{frame}{Esparsidade do modelo e dos dados}
\protect\phantomsection\label{esparsidade-do-modelo-e-dos-dados}
{\def\LTcaptype{none} % do not increment counter
\begin{longtable}[]{@{}
  >{\raggedright\arraybackslash}p{(\linewidth - 4\tabcolsep) * \real{0.3333}}
  >{\raggedright\arraybackslash}p{(\linewidth - 4\tabcolsep) * \real{0.3333}}
  >{\raggedright\arraybackslash}p{(\linewidth - 4\tabcolsep) * \real{0.3333}}@{}}
\toprule\noalign{}
\begin{minipage}[b]{\linewidth}\raggedright
Conceito
\end{minipage} & \begin{minipage}[b]{\linewidth}\raggedright
Significado
\end{minipage} & \begin{minipage}[b]{\linewidth}\raggedright
Exemplo
\end{minipage} \\
\midrule\noalign{}
\endhead
Modelo esparso & Poucos efeitos ativos & Cinco coeficientes não nulos \\
Matriz esparsa & Muitas entradas iguais a zero & Contagens de palavras
em documentos \\
\bottomrule\noalign{}
\end{longtable}
}

Uma matriz densa pode gerar um problema de regressão esparsa.

Uma matriz esparsa pode conter muitas variáveis relevantes.

\textbf{As duas propriedades não são equivalentes.}
\end{frame}

\begin{frame}{Lasso e seleção de variáveis}
\protect\phantomsection\label{lasso-e-seleuxe7uxe3o-de-variuxe1veis}
\textbf{Por que usar LASSO?} Quando muitas variáveis estão disponíveis,
mas esperamos poucos efeitos relevantes, podemos penalizar os
coeficientes para favorecer um modelo mais simples. A expressão abaixo
formaliza essa ideia.

O lasso estima

\[
\boxed{
\hat{\boldsymbol\beta}=
\arg\min_{\beta_0,\boldsymbol\beta}
\left\{
\frac{1}{2n}\sum_{i=1}^{n}
(y_i-\beta_0-\mathbf x_i^\top\boldsymbol\beta)^2
+\lambda\|\boldsymbol\beta\|_1
\right\}.
}
\]

A penalização \(L_1\) favorece coeficientes exatamente iguais a zero.

\begin{itemize}
\tightlist
\item
  \(\lambda\) controla a intensidade da regularização.
\item
  O intercepto fica fora da penalização.
\item
  A validação escolhe \(\lambda\) para o objetivo preditivo.
\end{itemize}
\end{frame}

\begin{frame}{Esparsidade também pode ser não linear}
\protect\phantomsection\label{esparsidade-tambuxe9m-pode-ser-nuxe3o-linear}
\textbf{Exemplo:} a resposta depende de poucas variáveis, mas por meio
de senos, quadrados e interações. A esparsidade existe na função, embora
um LASSO linear nas variáveis originais não consiga representar todas
essas relações.

\[
Y=\sin(X_2)+X_{17}^2+X_2X_{204}+\varepsilon.
\]

A função depende de poucas variáveis, mas inclui não linearidade e
interação.

\textbf{Lasso aplicado somente às entradas originais representa efeitos
lineares.}

Possibilidades:

\begin{itemize}
\tightlist
\item
  Bases e interações com penalização.
\item
  Modelos aditivos esparsos, quando a aditividade for adequada.
\item
  Métodos capazes de representar interações, com controle de
  complexidade.
\end{itemize}
\end{frame}

\section{Comparação das
estruturas}\label{comparauxe7uxe3o-das-estruturas}

\begin{frame}{Resumir ou selecionar?}
\protect\phantomsection\label{resumir-ou-selecionar}
{\def\LTcaptype{none} % do not increment counter
\begin{longtable}[]{@{}
  >{\raggedright\arraybackslash}p{(\linewidth - 4\tabcolsep) * \real{0.3333}}
  >{\raggedright\arraybackslash}p{(\linewidth - 4\tabcolsep) * \real{0.3333}}
  >{\raggedright\arraybackslash}p{(\linewidth - 4\tabcolsep) * \real{0.3333}}@{}}
\toprule\noalign{}
\begin{minipage}[b]{\linewidth}\raggedright
Aspecto
\end{minipage} & \begin{minipage}[b]{\linewidth}\raggedright
Redundância
\end{minipage} & \begin{minipage}[b]{\linewidth}\raggedright
Esparsidade
\end{minipage} \\
\midrule\noalign{}
\endhead
Estrutura & Informação compartilhada & Poucos efeitos relevantes \\
Operação natural & Construir representações & Selecionar variáveis ou
efeitos \\
Exemplo & 100 medidas de três fatores & Cinco entre 1.000 variáveis \\
Métodos associados & Ridge & LASSO \\
Risco & Perder sinal ao resumir & Seleção instável ou incorreta \\
\bottomrule\noalign{}
\end{longtable}
}

\textbf{As hipóteses podem coexistir.}

\note{Comparação e tamanho amostral estruturado: cerca de 15 minutos.}
\end{frame}

\begin{frame}{Esparsidade depende da representação}
\protect\phantomsection\label{esparsidade-depende-da-representauxe7uxe3o}
Considere entradas independentes e padronizadas e defina

\[
Z_1=\frac{X_1+\cdots+X_p}{\sqrt p},
\qquad Y=2Z_1+\varepsilon.
\]

Na representação original:

\[
\beta_j=\frac{2}{\sqrt p},\qquad j=1,\ldots,p.
\]

Todos os coeficientes são diferentes de zero, embora a resposta dependa
de \textbf{uma única combinação}.

\textbf{Dimensão preditiva pequena pode coexistir com coeficientes
densos.}
\end{frame}

\begin{frame}{Redundância e esparsidade juntas}
\protect\phantomsection\label{redunduxe2ncia-e-esparsidade-juntas}
Imagine 1.000 variáveis organizadas em 20 grupos correlacionados.

Apenas dois grupos contêm informação necessária à função de regressão.

\begin{itemize}
\tightlist
\item
  \textbf{Dentro dos grupos:} redundância.
\item
  \textbf{Entre os grupos:} esparsidade da estrutura relevante.
\end{itemize}

Possibilidades: elastic net, group lasso com grupos conhecidos ou
seleção seguida de representação.

Toda etapa aprendida dos dados deve integrar a validação.
\end{frame}

\begin{frame}{Alta dimensão na regressão linear}
\protect\phantomsection\label{alta-dimensuxe3o-na-regressuxe3o-linear}
\[
\mathbf y=X\boldsymbol\beta+\boldsymbol\varepsilon.
\]

Se \(p>n\), a matriz \(X^\top X\) é singular e os coeficientes de
mínimos quadrados não são únicos.

Para coeficientes penalizados e \(\lambda>0\), ridge utiliza

\[
\hat{\boldsymbol\beta}_{R}
=(X^\top X+\lambda I)^{-1}X^\top\mathbf y.
\]

A regularização torna a solução identificável pelo critério e controla
sua variabilidade.

\note{A fórmula assume entradas e resposta centradas, ou tratamento
separado do intercepto. Existência de uma solução única regularizada não
garante boa predição.}
\end{frame}

\begin{frame}{Tamanho amostral em modelos esparsos}
\protect\phantomsection\label{tamanho-amostral-em-modelos-esparsos}
Sob condições adequadas, limites de erro preditivo incluem termos da
ordem

\[
\boxed{\sigma^2\frac{s\log p}{n}.}
\]

{\def\LTcaptype{none} % do not increment counter
\begin{longtable}[]{@{}ll@{}}
\toprule\noalign{}
Termo & Papel \\
\midrule\noalign{}
\endhead
\(s\) & Quantidade de efeitos a estimar \\
\(\log p\) & Custo de procurar entre muitos candidatos \\
\(\sigma^2\) & Intensidade do ruído \\
\(n\) & Informação amostral disponível \\
\bottomrule\noalign{}
\end{longtable}
}

\textbf{Esparsidade reduz a complexidade, mas exige condições sobre os
preditores.}

\note{Limite típico para erro quadrático de predição de regressão linear
esparsa, com regularização calibrada, controle de caudas e condições de
compatibilidade ou autovalores restritos. Não é fórmula universal nem
garantia para suporte exato. Fonte: Gautier e Tsybakov, Regularity
Properties for Sparse Regression.}
\end{frame}

\begin{frame}{O que a expressão não fornece?}
\protect\phantomsection\label{o-que-a-expressuxe3o-nuxe3o-fornece}
A ordem \(s\log p\) não determina sozinha uma amostra recomendada.

A exigência real também depende de:

\begin{itemize}
\tightlist
\item
  Intensidade dos sinais e nível de ruído.
\item
  Correlação e geometria dos preditores.
\item
  Precisão desejada e objetivo da análise.
\item
  Capacidade de identificar a representação relevante.
\end{itemize}

\textbf{Prever bem, estimar coeficientes e recuperar o suporte são
objetivos distintos.}
\end{frame}

\begin{frame}{A dimensão relevante também exige aprendizagem}
\protect\phantomsection\label{a-dimensuxe3o-relevante-tambuxe9m-exige-aprendizagem}
Se uma representação útil \(\mathbf Z\in\mathbb R^d\) fosse conhecida,
poderíamos modelar

\[
Y=g(\mathbf Z)+\varepsilon.
\]

Quando precisamos estimá-la, aparece um custo adicional.

\begin{itemize}
\tightlist
\item
  Pouca amostra pode dificultar estimar componentes ou fatores.
\item
  A variância das entradas pode não acompanhar o sinal preditivo.
\item
  Uma representação de baixa dimensão não elimina ruído nem complexidade
  de \(g\).
\end{itemize}

\textbf{A redução de dimensão ajuda quando preserva informação útil e
pode ser aprendida.}
\end{frame}

\section{Escolha dos métodos}\label{escolha-dos-muxe9todos}

\begin{frame}{Árvores, random forest e boosting}
\protect\phantomsection\label{uxe1rvores-random-forest-e-boosting}
Árvores podem ignorar variáveis ao escolher divisões, mas muitas
candidatas irrelevantes dificultam encontrar sinal.

Na random forest, se há \(s\) variáveis relevantes entre \(p\), a
probabilidade de pelo menos uma delas estar entre \(q=\texttt{mtry}\)
candidatas é

\[
\boxed{
1-\frac{\binom{p-s}{q}}{\binom{p}{q}}.
}
\]

O boosting pode concentrar o ajuste em variáveis úteis, mas também
selecionar ruído ao longo das iterações.

\note{Fórmula combinatória para amostragem uniforme sem reposição das
candidatas. Convenção: binomial(p-s,q)=0 se q\textgreater p-s. Ter uma
variável relevante como candidata não garante que ela será escolhida nem
que a divisão será útil, sobretudo com interações. Poda/profundidade,
tamanho de folha, mtry, taxa de aprendizagem e iterações são controles
distintos.}
\end{frame}

\begin{frame}{Variáveis irrelevantes e candidatas na floresta}
\protect\phantomsection\label{variuxe1veis-irrelevantes-e-candidatas-na-floresta}
\pandocbounded{\includegraphics[keepaspectratio]{aula_outros_aspectos_files/figure-beamer/unnamed-chunk-5-1.png}}

Cálculo com \(s=5\) e \(\texttt{mtry}=10\) fixos. O gráfico ilustra
candidatas, sem medir erro preditivo.
\end{frame}

\begin{frame}{Redes neurais em altas dimensões}
\protect\phantomsection\label{redes-neurais-em-altas-dimensuxf5es}
A primeira camada totalmente conectada calcula

\[
\mathbf a^{(1)}=\phi\left(W^{(1)}\mathbf x+\mathbf b^{(1)}\right).
\]

Com \(H\) neurônios, ela contém

\[
\boxed{H(p+1)\text{ parâmetros}.}
\]

A rede pode aprender representações úteis, mas aumentar \(p\) também
aumenta os parâmetros.

\textbf{Arquitetura, regularização e informação amostral continuam
essenciais.}
\end{frame}

\begin{frame}{Quando uma RNA pode aproveitar estrutura?}
\protect\phantomsection\label{quando-uma-rna-pode-aproveitar-estrutura}
\begin{itemize}
\tightlist
\item
  Representações intermediárias podem combinar entradas redundantes.
\item
  Arquiteturas adequadas podem explorar composição e regularidades.
\item
  Weight decay e early stopping controlam o ajuste.
\item
  Uma camada com \(H<p\) cria uma representação menor, sem garantir
  preservar o sinal.
\end{itemize}

Em dados tabulares com \(p\gg n\), comparar a RNA com modelos
penalizados é essencial.

\textbf{Muitos parâmetros não fornecem, por si só, uma explicação para
boa generalização.}

\note{Resultados teóricos para redes explorando baixa dimensão
intrínseca exigem hipóteses sobre a distribuição e a classe de funções.
Evitar prometer que o treinamento vai encontrar automaticamente a
representação desejada. Fonte: Zhang et al.~(2023), Effective Minkowski
Dimension of Deep Nonparametric Regression.}
\end{frame}

\begin{frame}{Classificação em altas dimensões}
\protect\phantomsection\label{classificauxe7uxe3o-em-altas-dimensuxf5es}
\textbf{Exemplo:} em um diagnóstico baseado em milhares de medições, o
modelo precisa aprender uma regra de classificação e, depois,
transformar seus escores em decisões. Essas são duas tarefas distintas.

As hipóteses estruturais também orientam a classificação.

{\def\LTcaptype{none} % do not increment counter
\begin{longtable}[]{@{}ll@{}}
\toprule\noalign{}
Situação & Métodos a comparar \\
\midrule\noalign{}
\endhead
Efeitos aproximadamente lineares & Regressão logística penalizada \\
Separação por combinação linear & SVM linear regularizada \\
Fatores com informação preditiva & Componentes + classificador \\
Não linearidade e interações & Ensembles ou RNA regularizada \\
\bottomrule\noalign{}
\end{longtable}
}

Escolher métricas conforme o objetivo: log-loss, AUC, sensibilidade,
calibração ou custo de decisão.
\end{frame}

\section{Avaliação e discussão}\label{avaliauxe7uxe3o-e-discussuxe3o}

\begin{frame}{Validação sem vazamento de informação}
\protect\phantomsection\label{validauxe7uxe3o-sem-vazamento-de-informauxe7uxe3o}
\textbf{Exemplo:} selecionar variáveis usando toda a base e só depois
separar treino e teste faz o conjunto de teste influenciar a seleção. O
desempenho final tende a parecer melhor do que realmente é.

Em cada partição de treinamento da validação cruzada:

\begin{enumerate}
\tightlist
\item
  Estimar imputação e padronização.
\item
  Ajustar seleção, ACP ou outra representação.
\item
  Treinar o modelo com os hiperparâmetros candidatos.
\item
  Aplicar as transformações aprendidas à partição de validação.
\item
  Avaliar as predições nas observações reservadas.
\end{enumerate}

O conjunto de teste participa apenas da avaliação final.

\note{Encerramento e discussão: cerca de 10 minutos. Ajustar ACP no
banco inteiro também utiliza informação da distribuição das entradas de
validação/teste. Para estimar desempenho do pipeline de treinamento,
mantê-la dentro das partições. Para dados dependentes, usar divisões
temporais ou por grupos apropriadas.}
\end{frame}

\begin{frame}{Experimento para investigar as hipóteses}
\protect\phantomsection\label{experimento-para-investigar-as-hipuxf3teses}
{\def\LTcaptype{none} % do not increment counter
\begin{longtable}[]{@{}lll@{}}
\toprule\noalign{}
Cenário & Entradas & Função de regressão \\
\midrule\noalign{}
\endhead
Esparso & Preditores independentes & Cinco efeitos ativos \\
Redundante & Medições de três fatores & Depende dos três fatores \\
Misto & Grupos correlacionados & Depende de um grupo \\
\bottomrule\noalign{}
\end{longtable}
}

Comparar Ridge, LASSO e KNN.

Variar \(n\in\{100,300,1000\}\) e \(p\in\{20,100,500\}\).

\textbf{Manter controladas a variância do sinal e a variância do ruído
ao variar \(p\).}

\note{Este slide propõe um experimento, não apresenta resultados
empíricos. No cenário redundante, acrescentar sensores também pode
melhorar a medição dos fatores. Separar esse ganho da simples inclusão
de variáveis irrelevantes. Usar múltiplas repetições e conjuntos de
teste independentes. Um método oráculo com variáveis/fatores verdadeiros
pode servir como referência, identificado como tal.}
\end{frame}

\begin{frame}{O que avaliar no experimento?}
\protect\phantomsection\label{o-que-avaliar-no-experimento}
\textbf{Predição:} erro em teste e variabilidade entre repetições.

\textbf{Seleção:} quantas variáveis ativas foram recuperadas e quantas
irrelevantes entraram.

\textbf{Estabilidade:} quanto o modelo ou a seleção muda entre amostras.

\textbf{Custo computacional:} tempo de ajuste e escolha de
hiperparâmetros.

Um método pode prever bem mesmo quando a seleção de variáveis varia.
\end{frame}

\begin{frame}{Discussão em sala}
\protect\phantomsection\label{discussuxe3o-em-sala}
\begin{enumerate}
\tightlist
\item
  Acrescentar sensores redundantes equivale a acrescentar variáveis
  irrelevantes?
\item
  Quando ACP pode prejudicar o KNN ou a regressão?
\item
  Uma resposta que depende de uma combinação densa satisfaz esparsidade
  nas entradas originais?
\item
  Por que aumentar \(n\) pode ajudar sem tornar uma variável irrelevante
  informativa?
\item
  O melhor preditor precisa identificar exatamente as variáveis ativas?
\end{enumerate}

\note{Respostas esperadas: 1. Não. Sensores redundantes podem melhorar a
medição de fatores e alteram a geometria. 2. Quando descarta direções
preditivas ou mantém sobretudo variação irrelevante. 3. Não, embora
possa haver uma representação preditiva de dimensão baixa. 4. Aumenta
informação para estimar e validar, sem mudar a relevância populacional.
5. Não. Correlação pode permitir previsões semelhantes com suportes
diferentes.}
\end{frame}

\begin{frame}{Síntese}
\protect\phantomsection\label{suxedntese}
\textbf{Redundância:} muitas entradas compartilham informação. Podemos
buscar uma representação compacta.

\textbf{Esparsidade:} poucos efeitos são necessários. Podemos buscar
seleção de variáveis ou funções.

\textbf{Tamanho amostral:} depende da estrutura, do ruído e da precisão
desejada, além de \(p\).

\textbf{Escolha do método:} deve combinar hipóteses plausíveis e
avaliação adequada.

\[
\boxed{\text{Mais dados e estrutura adequada tornam a aprendizagem viável.}}
\]
\end{frame}

\section{Da previsão à decisão}\label{da-previsuxe3o-uxe0-decisuxe3o}

\begin{frame}{Por que uma boa previsão não basta?}
\protect\phantomsection\label{por-que-uma-boa-previsuxe3o-nuxe3o-basta}
\textbf{Exemplo: manutenção preventiva.} Um classificador estima que uma
máquina tem 30\% de probabilidade de falhar na próxima semana. Isso não
significa, por si só, que devemos interromper a produção.

Há duas perguntas diferentes:

\begin{enumerate}
\tightlist
\item
  \textbf{Previsão:} qual é a probabilidade de falha, dadas as medições?
\item
  \textbf{Decisão:} vale a pena parar a máquina agora, considerando os
  custos de cada erro?
\end{enumerate}

Em alta dimensão, a primeira pergunta já é difícil; a segunda exige
incorporar as consequências práticas da decisão.
\end{frame}

\begin{frame}{Probabilidade estimada e ponto de corte}
\protect\phantomsection\label{probabilidade-estimada-e-ponto-de-corte}
Suponha que o modelo produza
\(\hat\eta(\mathbf x)=P(Y=1\mid\mathbf X=\mathbf x)\), em que \(Y=1\)
representa falha.

Decidimos \textbf{intervir} quando \(\hat\eta(\mathbf x)>c\), para um
ponto de corte \(c\).

{\def\LTcaptype{none} % do not increment counter
\begin{longtable}[]{@{}
  >{\raggedright\arraybackslash}p{(\linewidth - 6\tabcolsep) * \real{0.2500}}
  >{\raggedleft\arraybackslash}p{(\linewidth - 6\tabcolsep) * \real{0.2500}}
  >{\raggedright\arraybackslash}p{(\linewidth - 6\tabcolsep) * \real{0.2500}}
  >{\raggedright\arraybackslash}p{(\linewidth - 6\tabcolsep) * \real{0.2500}}@{}}
\toprule\noalign{}
\begin{minipage}[b]{\linewidth}\raggedright
Máquina
\end{minipage} & \begin{minipage}[b]{\linewidth}\raggedleft
Probabilidade estimada
\end{minipage} & \begin{minipage}[b]{\linewidth}\raggedright
Decisão com \(c=0{,}50\)
\end{minipage} & \begin{minipage}[b]{\linewidth}\raggedright
Decisão com \(c=0{,}20\)
\end{minipage} \\
\midrule\noalign{}
\endhead
A & 0,12 & Não intervir & Não intervir \\
B & 0,30 & Não intervir & Intervir \\
C & 0,72 & Intervir & Intervir \\
\bottomrule\noalign{}
\end{longtable}
}

\textbf{Interpretação:} mudar o cutoff altera quais máquinas recebem
intervenção, \textbf{sem modificar as probabilidades estimadas}. Para
tomar decisões confiáveis, essas probabilidades devem estar
adequadamente calibradas.
\end{frame}

\begin{frame}{Por que o cutoff ótimo pode não ser 0,5?}
\protect\phantomsection\label{por-que-o-cutoff-uxf3timo-pode-nuxe3o-ser-05}
Imagine que uma intervenção desnecessária custe \textbf{R\$ 1.000} e
deixar passar uma falha custe \textbf{R\$ 9.000}.

Se \(\eta\) é a probabilidade de falha, os custos esperados são:

\[
\begin{aligned}
\text{Intervir: }&1000(1-\eta),\\
\text{Não intervir: }&9000\eta.
\end{aligned}
\]

Intervir é preferível quando \(1000(1-\eta)<9000\eta\), isto é,
\textbf{quando \(\eta>0{,}10\)}.

\textbf{Conclusão:} um limiar baixo pode ser racional quando perder uma
falha é muito mais caro que gerar um alarme falso. O cálculo pressupõe
custos fixos e probabilidades corretas.
\end{frame}

\begin{frame}{O erro de Bayes: por que nem o melhor modelo acerta tudo?}
\protect\phantomsection\label{o-erro-de-bayes-por-que-nem-o-melhor-modelo-acerta-tudo}
Considere dois equipamentos com medições quase idênticas. Um falha, o
outro não. Mesmo conhecendo perfeitamente a probabilidade condicional,
ainda existe incerteza.

Com perdas simétricas, a regra de Bayes escolhe a classe mais provável:

\[
g^*(\mathbf x)=I\{\eta(\mathbf x)>0{,}5\}.
\]

O menor erro condicional possível é
\(\min\{\eta(\mathbf x),1-\eta(\mathbf x)\}\).

\textbf{Exemplo:} se \(\eta(\mathbf x)=0{,}60\), escolher a classe 1 é
ótimo, mas ainda existe 40\% de chance de erro naquele perfil.
\textbf{Mais variáveis ou um algoritmo mais complexo não eliminam
necessariamente essa incerteza.}
\end{frame}

\section{Avaliação de
classificadores}\label{avaliauxe7uxe3o-de-classificadores}

\begin{frame}{Uma matriz de confusão conta uma história}
\protect\phantomsection\label{uma-matriz-de-confusuxe3o-conta-uma-histuxf3ria}
Em 1.000 máquinas, 50 realmente falham. Um modelo sinaliza 80, das quais
40 falham de fato.

{\def\LTcaptype{none} % do not increment counter
\begin{longtable}[]{@{}lrrr@{}}
\toprule\noalign{}
Situação real & Prever falha & Prever sem falha & Total \\
\midrule\noalign{}
\endhead
Falha & 40 (VP) & 10 (FN) & 50 \\
Sem falha & 40 (FP) & 910 (VN) & 950 \\
\bottomrule\noalign{}
\end{longtable}
}

O modelo acertou \(950/1000=95\%\) das decisões, mas gerou \textbf{40
alarmes falsos} e deixou passar \textbf{10 falhas}.

A acurácia, sozinha, não diz se o sistema é útil: precisamos saber
\textbf{quais erros importam}.
\end{frame}

\begin{frame}{Sensibilidade, especificidade e precisão}
\protect\phantomsection\label{sensibilidade-especificidade-e-precisuxe3o}
Para o exemplo anterior:

{\def\LTcaptype{none} % do not increment counter
\begin{longtable}[]{@{}
  >{\raggedright\arraybackslash}p{(\linewidth - 4\tabcolsep) * \real{0.3333}}
  >{\raggedright\arraybackslash}p{(\linewidth - 4\tabcolsep) * \real{0.3333}}
  >{\raggedright\arraybackslash}p{(\linewidth - 4\tabcolsep) * \real{0.3333}}@{}}
\toprule\noalign{}
\begin{minipage}[b]{\linewidth}\raggedright
Métrica
\end{minipage} & \begin{minipage}[b]{\linewidth}\raggedright
Cálculo
\end{minipage} & \begin{minipage}[b]{\linewidth}\raggedright
Leitura prática
\end{minipage} \\
\midrule\noalign{}
\endhead
Sensibilidade & \(40/50=80\%\) & Detecta 80\% das falhas reais \\
Especificidade & \(910/950\approx95{,}8\%\) & Descarta corretamente a
maioria das máquinas saudáveis \\
Precisão & \(40/80=50\%\) & Metade dos alertas corresponde a falhas
reais \\
\bottomrule\noalign{}
\end{longtable}
}

\textbf{Por que importa?} A equipe de manutenção recebe 80 alertas, mas
somente 40 são verdadeiros. Se inspecionar máquinas custa caro, a
precisão tem consequência operacional direta.
\end{frame}

\begin{frame}{Acurácia pode ser enganosa em classes raras}
\protect\phantomsection\label{acuruxe1cia-pode-ser-enganosa-em-classes-raras}
Agora suponha que apenas \textbf{1\%} das máquinas falhe. Um
classificador que sempre prevê ``sem falha'' terá \textbf{99\% de
acurácia}, embora não detecte nenhuma falha.

Por isso, em classes desbalanceadas, examine também
\textbf{sensibilidade, precisão, F1, balanced accuracy e custos
esperados}.

A \emph{balanced accuracy} é a média entre sensibilidade e
especificidade. Para o classificador que nunca detecta falhas, ela vale
\(0{,}5\), e não \(0{,}99\).

\textbf{Mensagem:} o desempenho precisa ser avaliado em relação ao
objetivo, não apenas pelo percentual global de acertos.
\end{frame}

\begin{frame}{Curva ROC: o que acontece ao mudar o cutoff?}
\protect\phantomsection\label{curva-roc-o-que-acontece-ao-mudar-o-cutoff}
Cada cutoff produz um par de resultados:

\begin{itemize}
\tightlist
\item
  \textbf{Taxa de verdadeiros positivos (TPR):} proporção de falhas
  detectadas.
\item
  \textbf{Taxa de falsos positivos (FPR):} proporção de máquinas
  saudáveis sinalizadas.
\end{itemize}

A curva ROC reúne esses pares à medida que variamos o cutoff. A
\textbf{AUC} resume a capacidade de ordenar positivos acima de
negativos: quanto maior, melhor a discriminação.

\textbf{Cuidado:} AUC não informa se as probabilidades estão calibradas
nem qual cutoff minimiza o custo. Dois modelos com AUC semelhante podem
gerar decisões bastante diferentes no limiar escolhido.
\end{frame}

\begin{frame}{Precision--Recall: quando o positivo é raro}
\protect\phantomsection\label{precisionrecall-quando-o-positivo-uxe9-raro}
Quando as falhas são raras, o interesse costuma ser: \textbf{entre os
alertas emitidos, quantos são verdadeiros?}

A curva Precision--Recall relaciona precisão e sensibilidade para
diferentes cutoffs.

\textbf{Exemplo:} detectar 90\% das falhas pode parecer excelente;
porém, se apenas 5\% dos alertas forem verdadeiros, o volume de
inspeções desnecessárias pode inviabilizar a operação.

A linha de base da precisão depende da prevalência de positivos. Por
isso, \textbf{não compare diretamente curvas PR obtidas em populações
com prevalências diferentes} sem contextualizar essa diferença.
\end{frame}

\begin{frame}{Três maneiras distintas de lidar com desbalanceamento}
\protect\phantomsection\label{truxeas-maneiras-distintas-de-lidar-com-desbalanceamento}
\textbf{1. Reamostragem:} altera a composição das observações usadas no
treinamento, por exemplo, aumentando a presença de casos raros.

\textbf{2. Pesos ou custos:} altera a importância dos erros durante o
ajuste do modelo.

\textbf{3. Cutoff:} mantém o modelo e altera a regra final de decisão.

Essas estratégias atuam em etapas diferentes e podem ser combinadas com
cuidado. \textbf{Reamostragem deve ocorrer somente dentro do
treinamento}, nunca antes da separação de teste; também pode afetar a
calibração das probabilidades.
\end{frame}

\begin{frame}{Combinar classificadores: por que isso pode ajudar?}
\protect\phantomsection\label{combinar-classificadores-por-que-isso-pode-ajudar}
Imagine três classificadores treinados com métodos ou amostras
diferentes. Eles podem cometer erros distintos.

\begin{itemize}
\tightlist
\item
  \textbf{Hard voting:} cada modelo vota em uma classe; prevalece a
  maioria.
\item
  \textbf{Soft voting:} combinamos probabilidades estimadas e depois
  aplicamos um cutoff.
\end{itemize}

\textbf{Exemplo:} probabilidades \(0{,}20\), \(0{,}65\) e \(0{,}75\)
resultam em média \(0{,}533\); com cutoff \(0{,}5\), o soft voting prevê
classe 1.

Combinar previsões pode reduzir instabilidade, mas \textbf{não garante
melhoria}: erros muito correlacionados e probabilidades mal calibradas
limitam os ganhos.
\end{frame}

\section{Interpretabilidade em alta
dimensão}\label{interpretabilidade-em-alta-dimensuxe3o}

\begin{frame}{Um modelo pode prever bem e ainda ser difícil de explicar}
\protect\phantomsection\label{um-modelo-pode-prever-bem-e-ainda-ser-difuxedcil-de-explicar}
Suponha que um modelo use \textbf{500 sensores} para prever falhas. A
equipe pergunta: ``A temperatura aumentou o risco ou apenas acompanhou a
pressão?''

Há pelo menos duas perguntas:

\begin{itemize}
\tightlist
\item
  \textbf{Global:} como as previsões do modelo costumam variar com a
  temperatura?
\item
  \textbf{Local:} por que este equipamento específico recebeu um alerta?
\end{itemize}

Interpretabilidade descreve o \textbf{comportamento do modelo}. Uma
explicação preditiva não demonstra, por si só, uma relação causal.
\end{frame}

\begin{frame}{PDP: efeito médio de uma variável na predição}
\protect\phantomsection\label{pdp-efeito-muxe9dio-de-uma-variuxe1vel-na-prediuxe7uxe3o}
O \emph{Partial Dependence Plot} (PDP) varia uma variável \(X_j\) e
calcula a média das previsões sobre as demais características
observadas:

\[
\widehat{PD}_j(z)=\frac{1}{n}\sum_{i=1}^{n}\hat f(z,\mathbf x_{i,-j}).
\]

\textbf{Como ler:} se a curva sobe com a temperatura, o modelo, em
média, produz previsões maiores ao substituir a temperatura por valores
crescentes.

\textbf{Limitação importante:} quando temperatura e pressão são
fortemente correlacionadas, a substituição pode criar combinações pouco
plausíveis, que não aparecem na operação real.
\end{frame}

\begin{frame}{ICE: a média pode esconder comportamentos diferentes}
\protect\phantomsection\label{ice-a-muxe9dia-pode-esconder-comportamentos-diferentes}
O PDP mostra uma \textbf{média}. Já as curvas \emph{Individual
Conditional Expectation} (ICE) mostram a resposta prevista para cada
observação quando variamos uma variável.

Imagine que aumentar a temperatura eleve o risco em equipamentos
antigos, mas quase não altere o risco em equipamentos novos. O PDP pode
suavizar essa diferença.

\textbf{Como interpretar:} curvas ICE aproximadamente paralelas sugerem
respostas individuais semelhantes; curvas com inclinações distintas
podem indicar interações ou heterogeneidade.

O PDP é a média das curvas ICE quando construídos sobre a mesma amostra
e grade de valores.
\end{frame}

\begin{frame}{Exemplo visual: PDP pode ocultar heterogeneidade}
\protect\phantomsection\label{exemplo-visual-pdp-pode-ocultar-heterogeneidade}
\pandocbounded{\includegraphics[keepaspectratio]{aula_outros_aspectos_files/figure-beamer/unnamed-chunk-6-1.png}}

\textbf{Leitura:} a curva média é constante, embora as respostas
individuais variem fortemente e em sentidos opostos. Este é um exemplo
construído para fins didáticos, não um resultado empírico.
\end{frame}

\begin{frame}{LIME: explicar uma previsão específica}
\protect\phantomsection\label{lime-explicar-uma-previsuxe3o-especuxedfica}
Uma máquina recebeu risco previsto de 82\%. \textbf{Quais
características levaram o modelo a produzir esse valor?}

O LIME gera perturbações ao redor dessa observação, consulta o modelo
original e ajusta um \textbf{modelo simples local} para aproximar suas
previsões.

A explicação pode indicar que, naquela vizinhança, temperatura alta
contribuiu positivamente e vibração baixa negativamente para a previsão
do modelo.

\textbf{Atenção:} contribuições LIME dependem de como a vizinhança e as
perturbações são definidas. Com variáveis correlacionadas, podem surgir
observações artificiais pouco realistas; a explicação pode variar entre
execuções.
\end{frame}

\begin{frame}{PDP, ICE e LIME: qual pergunta cada um responde?}
\protect\phantomsection\label{pdp-ice-e-lime-qual-pergunta-cada-um-responde}
{\def\LTcaptype{none} % do not increment counter
\begin{longtable}[]{@{}
  >{\raggedright\arraybackslash}p{(\linewidth - 4\tabcolsep) * \real{0.3333}}
  >{\raggedright\arraybackslash}p{(\linewidth - 4\tabcolsep) * \real{0.3333}}
  >{\raggedright\arraybackslash}p{(\linewidth - 4\tabcolsep) * \real{0.3333}}@{}}
\toprule\noalign{}
\begin{minipage}[b]{\linewidth}\raggedright
Ferramenta
\end{minipage} & \begin{minipage}[b]{\linewidth}\raggedright
Alcance
\end{minipage} & \begin{minipage}[b]{\linewidth}\raggedright
Pergunta respondida
\end{minipage} \\
\midrule\noalign{}
\endhead
PDP & Global & Como a predição média muda com \(X_j\)? \\
ICE & Individual, em conjunto & A relação é semelhante para todos? \\
LIME & Local & Como aproximar a previsão deste caso com uma explicação
simples? \\
\bottomrule\noalign{}
\end{longtable}
}

\textbf{Em alta dimensão, interpretabilidade exige atenção à
redundância:} se muitos preditores carregam a mesma informação, atribuir
importância a uma variável isolada pode ser instável ou enganoso.
\end{frame}

\begin{frame}{Um fluxo completo: do treinamento à decisão}
\protect\phantomsection\label{um-fluxo-completo-do-treinamento-uxe0-decisuxe3o}
\textbf{Exemplo: manutenção preditiva com 500 sensores e falhas raras.}

\begin{enumerate}
\tightlist
\item
  Separar treinamento, validação e teste respeitando a estrutura
  temporal, quando aplicável.
\item
  Comparar modelos regularizados e verificar o efeito da redundância e
  da esparsidade.
\item
  Avaliar discriminação, calibração e métricas adequadas à raridade das
  falhas.
\item
  Escolher o cutoff conforme os custos de alarmes falsos e falhas não
  detectadas.
\item
  Interpretar previsões com PDP, ICE ou LIME, verificando combinações
  irreais de sensores.
\item
  Avaliar uma única vez no teste final, sem usá-lo para refinar o
  procedimento.
\end{enumerate}

\textbf{Conclusão:} prever, decidir, avaliar e interpretar são etapas
relacionadas, mas não intercambiáveis.
\end{frame}

\begin{frame}{Para verificar sua compreensão}
\protect\phantomsection\label{para-verificar-sua-compreensuxe3o}
\textbf{Questão 1.} Por que um KNN pode piorar quando acrescentamos 100
variáveis irrelevantes, mesmo após padronização?

\textbf{Questão 2.} Um classificador com 99\% de acurácia pode ser
inútil para detectar falhas raras? Explique.

\textbf{Questão 3.} Se o custo de não detectar uma falha aumenta, o
cutoff ótimo tende a subir ou descer, mantendo os demais custos fixos?

\textbf{Questão 4.} Por que um PDP pode ser enganoso quando duas
variáveis são fortemente correlacionadas?

\note{Respostas: (1) distâncias incorporam ruído e deixam de refletir
proximidade relevante; (2) sim, prever sempre a classe majoritária pode
gerar alta acurácia e sensibilidade zero; (3) descer; (4) substituições
marginais criam combinações pouco plausíveis.}
\end{frame}

\begin{frame}{Referências para aprofundamento}
\protect\phantomsection\label{referuxeancias-para-aprofundamento}
\begin{itemize}
\tightlist
\item
  James, Witten, Hastie e Tibshirani (2021). \emph{An Introduction to
  Statistical Learning}, 2ª ed.~Capítulos 6, 7 e 10.
  \href{https://www.statlearning.com/}{statlearning.com}.
\item
  Hastie, Tibshirani e Friedman (2009). \emph{The Elements of
  Statistical Learning}, 2ª ed.~Capítulos 2, 3 e 18.
  \href{https://hastie.su.domains/ElemStatLearn/}{Livro}.
\item
  Zou e Hastie (2005). \emph{Regularization and variable selection via
  the elastic net}. JRSS B, 67, 301--320.
  \href{https://web.stanford.edu/~hastie/Papers/elasticnet.pdf}{Artigo}.
\end{itemize}
\end{frame}

\begin{frame}{Referências dos resultados teóricos}
\protect\phantomsection\label{referuxeancias-dos-resultados-teuxf3ricos}
\begin{itemize}
\tightlist
\item
  Wasserman. \emph{Nonparametric Regression}. Notas de Statistical
  Machine Learning, Carnegie Mellon.
  \href{https://stat.cmu.edu/~larry/=sml/nonpar.pdf}{Notas}.
\item
  Gautier e Tsybakov (2016). \emph{Regularity Properties for Sparse
  Regression}.
  \href{https://pmc.ncbi.nlm.nih.gov/articles/mid/NIHMS750383/}{Artigo}.
\item
  Zhang et al.~(2023). \emph{Effective Minkowski Dimension of Deep
  Nonparametric Regression: Function Approximation and Statistical
  Theories}. ICML, PMLR 202, 40911--40931.
  \href{https://proceedings.mlr.press/v202/zhang23f.html}{Artigo}.
\item
  Stanford, STATS 202. \emph{Dimensionality reduction}.
  \href{https://web.stanford.edu/class/stats202/notes/Model-selection/Dimensionality-reduction.html}{Notas}.
\end{itemize}
\end{frame}




\end{document}
